% =====================================================================
% CAPITOLO 7 — ORION: VALUTAZIONE SPERIMENTALE
% =====================================================================
% POLPA, parte B3 di B. E' il capitolo piu' originale e piu' "scientifico"
% della tesi: il framework di valutazione iterativa del system prompt
% del router e' uno dei tre contributi principali.
%
% Riferimento: plans/ORION_ROADMAP.md (sez. 6).
% Ordine di scrittura: TERZO (dopo Cap. 3-4-5-6), parallelamente alla
% raccolta dei risultati.
% =====================================================================

\chapter{Orion: valutazione sperimentale}
\label{chap:orion-valutazione}

% =====================================================================
\section{Domande di ricerca e impianto sperimentale}
\label{sec:eval-research-questions}
% =====================================================================
% TODO — Esplicitare le domande di ricerca a cui questo capitolo
% risponde:
%
% RQ1 — Quanto e' accurato il routing dell'orchestratore (agent accuracy)
%       sul dominio Coop?
% RQ2 — Quanto e' accurata la stima di complessita' (complexity accuracy)?
% RQ3 — Il confidence score e' calibrato? (E' piu' alto sui casi
%       corretti che su quelli errati?)
% RQ4 — Il meccanismo di clarification riduce l'errore senza degradare
%       l'esperienza utente (clarification rate non eccessivo)?
% RQ5 — Iterazioni del system prompt migliorano le metriche in modo
%       misurabile?
%
% Citare \cite{liu2023agentbench} come framework di riferimento per la
% valutazione di agenti LLM e \cite{zheng2023llmjudge} per la
% metodologia LLM-as-judge (eventualmente usata per valutare il
% reasoning del router come metrica complementare).

% =====================================================================
\section{Metodologia di valutazione}
\label{sec:eval-metodologia}
% =====================================================================

\subsection{Costruzione del gold standard su Coda}
\label{subsec:eval-gold-standard}
% TODO — Test suite Coda: tabella "Test Suite Orion" con colonne
% ID, Prompt, Agente Atteso, Complessita' Attesa, Categoria, Note,
% Attivo. Numero minimo di test case: 50 (10 per agente + 10 ambigui).
%
% Costruzione: prompt sintetizzati a partire da log reali della suite
% (anonimizzati) + prompt creati ad hoc per coprire i casi limite
% (intent overlap tra SysAid e AllyCare; richieste fuori dominio;
% domande ambigue tra Vera e CooPolicy).
%
% Categorie di prompt: SysAid, NPS, Policy, General, Ambiguous.

% [FIGURA 7.1] Schema della tabella Coda con un estratto rappresentativo
% dei test case (5-7 righe).
%
% \begin{figure}[h]
%     \centering
%     \includegraphics[width=\textwidth]{fig07_coda_test_suite.pdf}
%     \caption[Schema della test suite su Coda]{Estratto della tabella \texttt{Test Suite Orion} su Coda: gold standard di prompt etichettati con agente atteso, complessita' attesa e categoria tematica.}
%     \label{fig:coda_test_suite}
% \end{figure}

\subsection{Tabella dei risultati e metriche derivate}
\label{subsec:eval-results-table}
% TODO — Tabella "Risultati Valutazione Orion": una riga per esecuzione
% di test case (Run ID + Test Case ID + Versione System Prompt come
% chiave composta). Colonne: agente atteso vs prodotto, complessita'
% attesa vs prodotta, confidence, reasoning, formula di accuracy.

\subsection{Lo script di valutazione (\texttt{orion-eval/})}
\label{subsec:eval-script}
% TODO — Architettura dello script:
% orion-eval/
%   main.py           # entry point
%   coda_client.py    # wrapper Coda API
%   orion_client.py   # chiama POST /api/orion/classify
%   report.py         # report testuale + CSV
%   config.py         # env vars
%
% Flusso: legge i test case attivi, chiama /classify, confronta con
% atteso, scrive su Coda, stampa summary.

% =====================================================================
\section{Metriche e target}
\label{sec:eval-metriche}
% =====================================================================
% TODO — Tabella delle metriche con target dichiarati a priori:
%
% Agent Accuracy:        correct_agent / total          target >= 85%
% Complexity Accuracy:   correct_complexity / total     target >= 80%
% Clarification Rate:    clarification_asked / total    target < 15%
% Mean Confidence (cor.): media confidence sui corretti target >= 0.82
% Mean Confidence (err.): media confidence sugli errati  target <= 0.65
%
% L'ultima coppia di metriche valida la CALIBRAZIONE del confidence
% score: un buon router deve essere piu' incerto quando sbaglia
% \cite{kadavath2022knowsknows}.

% =====================================================================
\section{Risultati}
\label{sec:eval-risultati}
% =====================================================================
% TODO — Da popolare con i numeri reali al termine degli esperimenti.
% Strutturare i risultati in 4 sotto-sezioni, una per gruppo di metriche.
%
% IMPORTANTE — Per ogni risultato presentato:
% 1. Domanda di ricerca a cui risponde (RQ1..5).
% 2. Numeri concreti (con unita' di misura).
% 3. Confronto con il target.
% 4. Discussione dei casi anomali (se presenti).

\subsection{Routing accuracy per agente}
\label{subsec:eval-routing-accuracy}
% TODO — Tabella + bar chart: accuracy per ciascun agente target
% (self, sysaid_gemini, sysaid_sonnet, allycare, coopolicy).
% Identificare gli agenti su cui l'accuracy e' piu' bassa e ipotizzare
% le ragioni.

% [FIGURA 7.2] Bar chart accuracy per agente. DATO REALE da popolare.
%
% \begin{figure}[h]
%     \centering
%     \includegraphics[width=0.85\textwidth]{fig07_accuracy_per_agent.pdf}
%     \caption[Accuracy di routing per agente target]{Accuracy di routing dell'orchestratore Orion, suddivisa per agente target. La linea tratteggiata indica il target di progetto (85\%).}
%     \label{fig:accuracy_per_agent}
% \end{figure}

\subsection{Calibrazione del confidence score}
\label{subsec:eval-confidence-cal}
% TODO — Confronto tra distribuzione di confidence sui casi corretti
% e quella sui casi errati. Analisi: il modello e' calibrato? La soglia
% 0.60 e' la migliore?
% Strumenti possibili: reliability diagram, expected calibration error.

% [FIGURA 7.3] Istogramma confidence: corretti vs errati.
% DATO REALE da popolare.
%
% \begin{figure}[h]
%     \centering
%     \includegraphics[width=0.85\textwidth]{fig07_confidence_calibration.pdf}
%     \caption[Calibrazione del confidence score]{Distribuzione del confidence score sui casi correttamente classificati e su quelli erroneamente classificati. Una buona calibrazione richiede che le due distribuzioni siano separate.}
%     \label{fig:confidence_calibration}
% \end{figure}

\subsection{Iterazione del system prompt}
\label{subsec:eval-iterazione}
% TODO — Documentare le iterazioni: v1.0 -> v1.1 -> v1.2 -> ...
% Per ogni iterazione: cosa e' stato cambiato nel system prompt, cosa
% e' migliorato, cosa e' peggiorato. Trade-off osservati.
% Citare \cite{sahoo2024prompteng_survey} come quadro teorico del
% prompt engineering iterativo.

% [FIGURA 7.4] Timeline accuracy per versione di system prompt.
% DATO REALE.
%
% \begin{figure}[h]
%     \centering
%     \includegraphics[width=0.9\textwidth]{fig07_prompt_iteration.pdf}
%     \caption[Iterazione del system prompt]{Andamento delle metriche di accuracy lungo le versioni successive del system prompt del router.}
%     \label{fig:prompt_iteration}
% \end{figure}

\subsection{Latenza del routing step}
\label{subsec:eval-latenza}
% TODO — Misurazione della latenza p50/p95/p99 della call al router.
% Importante per valutare l'impatto sull'esperienza utente: il routing
% step aggiunge latenza all'intera pipeline.
% Misurare anche la latenza del proxy (overhead Cloud Run-to-Cloud Run).

% [FIGURA 7.5] Distribuzione latenza routing step (boxplot p50/p95/p99).
%
% \begin{figure}[h]
%     \centering
%     \includegraphics[width=0.8\textwidth]{fig07_latenza.pdf}
%     \caption[Latenza del routing step]{Distribuzione della latenza del routing step (boxplot p50/p95/p99) misurata su un campione rappresentativo di richieste.}
%     \label{fig:routing_latency}
% \end{figure}

% =====================================================================
\section{Analisi dei casi di errore}
\label{sec:eval-errori}
% =====================================================================
% TODO — Analisi qualitativa dei casi in cui Orion sbaglia. Pattern
% ricorrenti identificati nei reasoning del router. Esempi di
% "false positive" e "false negative" rispetto a ciascun agente.
% Lezioni apprese che hanno portato alle iterazioni del system prompt.

% =====================================================================
\section{Discussione: limiti dell'impianto di valutazione}
\label{sec:eval-limiti}
% =====================================================================
% TODO — Discutere onestamente i limiti:
% - Dimensione del gold standard (50 test case e' un minimo).
% - Bias del costruttore del dataset (l'autore ha costruito i prompt).
% - Dataset sintetico vs traffico reale: gap potenziale.
% - LLM-as-judge per il reasoning: noto che ha bias \cite{zheng2023llmjudge}.
% - Confronto con \cite{liang2022helm} (HELM): la valutazione qui e'
%   piu' specifica del contesto, meno "olistica".
